\section{The Map and Capacity (The Dirac Propagator)}
\label{sec:map}

\textbf{The Architectural Goal:} 
\begin{itemize}
    \item \textbf{Map ($\bar{\psi} i \gamma^\mu D_\mu \psi$):} \textit{The Cost of Propagation.} The Dirac operator is the logic of state evolution. The spinor structure ($\gamma^\mu$) is the algebraic requirement for a complex signal to propagate on the chiral $D_4$ substrate.
    \item \textbf{Capacity ($-m\bar{\psi}\psi$):} \textit{The Cost of Existence.} The mass term is the impedance required to lock a persistent state into a stable history against the vacuum flux. It represents the phase-locked memory of the lattice.
\end{itemize}

\textbf{Terminology Mapping:} Throughout this text, we standardize the terminology bridging geometry and physics. A localized, persistent boundary structure on the substrate is a \textbf{Topological Defect}. When this defect occupies the chiral node channels ($\nu=16$), its physical manifestation is a \textbf{Fermion}. 

We first derive the dynamics of the \textbf{Map} pillar by defining the evolution of a fermion state $\psi(x,t)$ on the lattice. Time is treated not as a continuous dimension, but as a discrete clock cycle defined by the fundamental temporal step $\tau$.

\subsection{The Unitary Flux Constraint}
The evolution of the state is governed by the conservation of probability amplitude (Unitarity). We define the update rule based on the \textbf{Quantum Lattice Gas Automaton (QLGA)} formalism \cite{meyer_quantum_1996}. The QLGA is the uniquely appropriate discretization scheme for this substrate because it explicitly enforces the three core requirements derived for the vacuum in System 0: locality (nearest-neighbor transitions), unitarity (probability conservation at each discrete step), and finiteness (a strictly bounded state space per node). It generalizes the Feynman Checkerboard model \cite{feynman_quantum_1965} to the $D_4$ geometry.

The state at a node $\psi(x, t+\tau)$ is the superposition of its previous state and the information flux received from its chiral neighbors. This is the \textbf{Kirchhoff Condition for Information}: the net flux into a node determines the rate of change of the state amplitude.

We define the discrete evolution operator as:
\begin{align}
    \psi(x, t + \tau) = &\psi(x, t) \nonumber \\
    &- ig \sum_{k=1}^3[\psi(x + e_k \ell, t) - \psi(x - e_k \ell, t)] \nonumber \\
    &- im\psi(x, t)
\label{eq:unitary_flux}
\end{align}
Where:
\begin{itemize}
    \item $e_k$ are the spatial lattice basis vectors.
    \item $\ell$ is the spatial node spacing.
    \item $g$ is the vacuum admittance (coupling efficiency between nodes).
    \item $m$ is the effective topological impedance ($m = y\langle\phi\rangle$), representing the geometric cost of state retention mediated by the scalar background.
    \item The factor $i$ ensures unitary rotation rather than diffusive decay.
\end{itemize}

\subsection{The Continuum Limit}
To recover the field equation, we rearrange the update rule into difference quotients. We take the limit where the lattice spacing $\ell$ and time step $\tau$ approach zero. Introducing the required factor of 2 for the central difference span, we pull out the causal bandwidth ratio $\ell / \tau$:

\begin{equation}
\begin{split}
    \frac{\psi(t+\tau) - \psi(t)}{\tau} =& \\ 
    -i g \sum_{k=1}^{3} \left( \frac{\psi(x + e_k \ell) - \psi(x - e_k \ell)}{2\ell} \right) &\frac{2\ell}{\tau} \\
    - i m \psi(t)&
    \end{split}
\end{equation}

In the continuum limit, we enforce the kinematic capacity constraint $c \equiv \ell / \tau = 1$ and absorb the normalized geometric admittance ($2gc \to 1$). Crucially, $g$ is not a tunable coupling parameter as treated in standard effective field theories. It is the intrinsic geometric transmission efficiency of the lattice, strictly forced to $g = 1/2c$ by the requirement of lossless unitary propagation at the causal bandwidth limit. This is the first instance of the framework's overarching geometric constraint: the substrate leaves no continuous free parameters. 

Recognizing the central difference quotient as the spatial derivative $\partial_k$ and the left-hand side as the temporal derivative $\partial_t$:

\begin{equation}
    \partial_t \psi = -i \sum_{k=1}^{3} \Gamma^k \partial_k \psi - i m \psi
\end{equation}

Multiplying by $i$ to match standard quantum mechanical notation yields the Hamiltonian form:
\begin{equation}
    i \partial_t \psi = \sum_{k=1}^{3} \Gamma^k \partial_k \psi + m \psi
\end{equation}
Here, $\Gamma^k$ represents the projection of the lattice vectors onto the propagation manifold.

\subsection{The Clifford Identification}
The critical step in identifying this lattice rule with the Dirac equation is the geometric nature of the propagation matrices $\Gamma^k$. In the $E_8 \to D_4 \oplus D_4$ projection, the lattice basis vectors are not simple coordinate directions; they are \textbf{Spinor Generators}.

For the propagator to be consistent on the projection manifold, the full set of spacetime basis vectors $e_\mu = (e_0, e_k)$ must satisfy the metric of that manifold. The anti-commutation relation for the $D_4$ causal lattice basis is:
\begin{equation}
    \{e_\mu, e_\nu\} \equiv e_\mu e_\nu + e_\nu e_\mu = 2\eta_{\mu\nu} \mathbb{I}
\end{equation}
where $\eta_{\mu\nu} = \text{diag}(-1,+1,+1,+1)$ is the Lorentzian metric signature derived from the causal projection.

This relation is precisely the defining algebra of the Dirac gamma matrices ($\gamma^\mu$). Therefore, we explicitly identify the lattice update vectors $e_\mu$ with the gamma matrices:
\begin{equation}
    (e_0, e_1, e_2, e_3) \equiv (\gamma^0, \gamma^1, \gamma^2, \gamma^3)
\end{equation}

To recover the covariant Dirac equation from the Hamiltonian form, we recognize that the spatial QLGA matrices correspond to the Dirac $\alpha$ matrices ($\Gamma^k = \gamma^0 \gamma^k$), and the mass impedance structure aligns with the $\beta$ matrix ($\gamma^0$). Multiplying the entire continuum equation by the temporal vector $\gamma^0$ perfectly recovers the covariant Dirac equation:
\begin{equation}
    (i\gamma^\mu \partial_\mu - m)\psi = 0
\end{equation}

\subsection{The Chiral Filter: Geometric Parity Violation}
Standard QFT introduces the chiral projection operator $P_L = \frac{1-\gamma^5}{2}$ as an empirical constraint to match the observed weak interaction. In the \textit{$E_8$-Persistence} framework, this is a geometric necessity of the projection.

The decomposition $E_8 \to D_4 \oplus D_4$ splits the 16-component lattice state into two orthogonal sectors. Crucially, the causal metric signature $(-1, +1, +1, +1)$ required for the \textbf{Toll} pillar (Time) can only be assigned to one $D_4$ sector (Spacetime). The orthogonal $D_4$ sector (Internal Symmetry) remains Euclidean $(+1, +1, +1, +1)$.

Dynamics requires a time derivative ($\partial_t$). States residing in the Euclidean sector possess no temporal generator; they are geometrically static ($\partial_t \psi_{mirror} \equiv 0$) relative to the causal manifold. Consequently, the ``Right-Handed'' (Mirror) components of the $E_8$ spinor are structurally filtered out of the propagation dynamics.
The Dirac propagator on the projected manifold acts only on the active sub-algebra, enforcing maximum parity violation ($P_L$) not by suppression, but by the topological absence of a temporal update path for the mirror sector.

\subsection{The Capacity Mechanism: Mass as Chiral Impedance}

If the vacuum only possessed the Map pillar, the resulting massless Dirac equation ($\mathcal{L}_{Map} \propto \bar{\psi} i\gamma^\mu D_\mu \psi$) would describe pure ballistic transport at $v=c$. While this perfectly satisfies information propagation, it completely violates the \textbf{Capacity} pillar. A signal moving at the speed of light possesses no rest frame; it is entirely delocalized along its light cone and thus cannot store local configuration memory.

\textbf{Geometric Necessity of the Scalar Bridge:} To satisfy Capacity, the node must retain state over time. In the lattice formalism, maintaining a local address requires creating a resonant geometric loop: converting an outgoing flux vector (Left-Chiral) into an incoming flux vector (Right-Chiral), effectively bouncing the signal back and forth to arrest its spatial translation (a geometric ``Zitterbewegung'').

However, as established by the Chiral Filter, the Left and Right chiral sectors reside in strictly orthogonal $D_4$ projections. Consequently, they cannot couple directly. A bare mass term ($-m\bar{\psi}\psi = -m(\bar{\psi}_L\psi_R + \bar{\psi}_R\psi_L)$) represents a direct transition between orthogonal spaces, which is geometrically forbidden. This structural prohibition recovers the chiral gauge invariance of the Standard Model not as an abstract gauge symmetry, but as a rigid geometric absolute.

\textbf{The Yukawa Operator:} To arrest information flow and satisfy Capacity without violating this geometric isolation, the lattice must introduce a mediating scalar variable—the \textbf{Governor} field ($\phi$), the architectural mechanism preventing unbounded divergence \cite{meyer_informational_2026}—to transfer momentum and phase between the orthogonal sectors. The Capacity constraint therefore uniquely mandates the Yukawa Coupling:
\begin{equation}
    \mathcal{L}_{Cap} \propto y \left( \bar{\psi}_L \phi \psi_R + \bar{\psi}_R \phi^\dagger \psi_L \right)
\end{equation}

\textbf{Conclusion: The Physical Nature of Mass.} This derivation reframes the physical nature of mass. The Dirac equation is not an axiom, nor are its chiral gauge restrictions arbitrary; it is the hydrodynamic limit of a discrete, unitary cellular automaton evolving on a spinor lattice. The ``matrices'' of the Standard Model are simply the basis vectors of the underlying geometry. Crucially, fermion mass is not an intrinsic, fundamental property of a particle. It is the geometric impedance ($m = y\langle\phi\rangle$) generated by the scalar boundary field—the necessary thermodynamic drag required to arrest a state from dissipating at the speed of light, locking a topological defect into a persistent local capacity.